\section*{Nomenclatura}
\addcontentsline{toc}{section}{Nomenclatura}
\label{sec:nomenclatura}

Los símbolos se agrupan en cinco bloques (índices y conjuntos, variables de estado,
parámetros del sistema, funciones y operadores, y métricas de evaluación). Entre
paréntesis o corchetes se indica el dominio o la unidad cuando procede.

\begingroup
\linespread{1.2}\selectfont
\setlength{\tabcolsep}{3pt}
\renewcommand{\arraystretch}{1.15}
\setlength{\LTpre}{6pt}
\setlength{\LTpost}{6pt}

% Cabecera de bloque: etiqueta naranja que ocupa el ancho de la tabla.
\newcommand{\nomhead}[1]{\addlinespace[2pt]%
  \multicolumn{2}{@{}l}{\bfseries\color{VIUOrange}#1}\\*[1pt]}

\begin{longtable}{@{}l@{\hskip 2em}>{\raggedright\arraybackslash}p{0.72\linewidth}@{}}
\toprule
\nomhead{Índices y conjuntos}
$i$                & Índice de robot ($\{1,\ldots,N\}$) \\
$k$                & Índice de tarea; $k{=}0$ representa la inactividad ($\{0,\ldots,K\}$) \\
$t$                & Paso temporal discreto ($\mathbb{N}$) \\
$\mathcal{S}$      & Conjunto de estrategias ($\{0,\ldots,K\}$) \\
$\mathcal{N}_i(t)$ & Vecinos de $i$ en el grafo de comunicación \\
$\mathcal{T}_i(t)$ & Tareas visibles o conocidas localmente por el robot $i$ \\
$\mathcal{O}_i(t)$ & Obstáculos detectados por el robot $i$ dentro de su radio de percepción \\
$\mathcal{C}_k(t)$ & Coalición formada: asignados a $k$ y presentes en $r_{\mathrm{det}}$ \\
\midrule
\nomhead{Variables de estado}
$\mathbf{q}_i$         & Pose del robot $i$: $(x,y,\theta)$ [m, m, rad] \\
$\mathbf{q}_i^{xy}$    & Proyección plana de la pose [m] \\
$p_{ik}$               & Probabilidad mixta del robot $i$ sobre la tarea $k$ ($[0,1]$) \\
$y_{ik}$               & Preferencia mixta usada en el desarrollo teórico avanzado; equivalente continuo de $p_{ik}$ \\
$\hat{x}_{ik}$         & Fracción estimada de robots en la tarea $k$ ($[0,1]$) \\
$\hat{C}_{ik}$         & Capacidad acumulada estimada por $i$ para la tarea $k$ [kg] \\
$s_{ik}$               & Capacidad efectiva de $i$ para $k$, $c_i\Psi(T_k(p_i))$ \\
$\zeff_k$              & Ocupación efectiva descontada espacialmente \\
$\Zeff_k$              & Capacidad efectiva agregada de la tarea $k$ \\
$e^Z_{ik}$             & Error local del DAC para el agregado de capacidad de la tarea $k$ \\
$s_{ik}^{W}$           & Contribución útil del robot $i$ en espacio wrench para la tarea $k$ \\
$\delta_k^W$           & Déficit de wrench útil de la tarea $k$ \\
$\gamma_i$             & Fase continua del robot $i$ sobre el contorno suavizado de la carga \\
$p_i^\star$            & Punto cartesiano objetivo asociado a la fase de contacto $\gamma_i$ [m] \\
$\mathcal E_W$         & Residual físico de wrench entre demanda de la carga y aporte de la coalición \\
$X^N(t)$               & Ocupación empírica pura de una población finita de $N$ robots \\
$a_{ik},b_{ik}$        & Factores genéricos de agregado y payoff usados en la condición de integrabilidad \\
$a_k^p$                & Masa de robots de clase $p$ asignada a la tarea $k$ \\
$Q$                    & Capacidad total desplegada por la flota \\
$C_p,A_p$              & Capacidad de la clase $p$ y capacidad acumulada desde clases pesadas \\
$h_i$                  & Margen de retornabilidad energética del robot $i$ \\
$\tilde{x}_k$          & Ocupación empírica pura, $a_k/N$ ($[0,1]$) \\
$a_k$                  & Ocupación asignada pura, $\sum_i\mathbf{1}_{[s_i=k]}$ \\
$\mathrm{count}_k$     & Tamaño de la coalición, $|\mathcal{C}_k|$; cumple $\mathrm{count}_k \le a_k$ \\
$s_i$                  & Estrategia pura actual del robot $i$ \\
$\mathbf{h}_i$         & Punto adelantado (\emph{look-ahead}) de navegación [m] \\
$\bar{\mathbf{h}}_k$   & Centroide de los puntos adelantados de la coalición $\mathcal{C}_k$ [m] \\
$\mathrm{cnt}^{\pm}_k$ & Contadores de pasos para la transición de modo \\
\midrule
\nomhead{Parámetros del sistema}
$N,K$                   & Número de robots y de tareas \\
$n_k$                   & Cardinalidad mínima exigida por la tarea $k$ \\
$T_s$                   & Período de muestreo del bucle de control [s] \\
$\beta$                 & Pendiente de la sigmoide de payoff \\
$\tilde\beta$           & Pendiente reescalada para el límite fluido, $\beta_N=\tilde\beta/N$ \\
$\beta_c$               & Pendiente de la sigmoide de capacidad heterogénea \\
$\nu_k$                 & Demanda fraccional de la tarea $k$, $n_k^{(N)}=\nu_kN$ \\
$\lambda$               & Escala espacial del payoff [m] \\
$r_k$                   & Recompensa base de la tarea $k$ \\
$r_0$                   & Recompensa de inactividad y precio de reserva del robot marginal \\
$\rho_p$                & Precio de reserva por unidad de capacidad de la clase $p$, $\rho_0/c^p$ \\
$\mu_0,\mu_{\min}$      & Tasa de revisión base y piso de la tasa \\
$\sigma_{\mathrm{rev}}$ & Ancho espacial de la tasa de revisión [m] \\
$\tau_s$                & Pasos necesarios para activar el transporte \\
$\rho_s$                & Umbral de histéresis de la transición de modo \\
$\ell$                  & Distancia de \emph{look-ahead} [m] \\
$\Delta T_c$            & Período de replanificación del tratamiento T5 [s] \\
$\varepsilon$           & Ganancia de consenso \\
$\varepsilon_r$         & Regularización del término de repulsión [m] \\
$\varepsilon_f$         & Regularización del término de formación [m] \\
$r_{\mathrm{det}}$      & Radio de detección de coalición [m] \\
$r_{\mathrm{form}}$     & Radio de activación de la formación [m] \\
$R$                     & Rango de comunicación [m] \\
$R_c$                   & Radio local de comunicación y vecindad para rigidez/eventos [m] \\
$R_{\mathrm{vis}}$      & Radio de percepción local de tareas, robots y obstáculos [m] \\
$\rho_f$                & Radio del anillo de formación, $r_{\mathrm{load}}+r_{\mathrm{robot}}$ [m] \\
$d^*$                   & Separación adyacente emergente en la formación, $2\rho_f\sin(\pi/n_k)$ [m] \\
$\alpha$                & Ganancia de atracción del potencial \\
$\eta$                  & Ganancia de repulsión inter-robot \\
$\delta$                & Ganancia de formación \\
$v_{\max},\omega_{\max}$ & Velocidad lineal y angular máximas [m/s, rad/s] \\
$\omega_{R,i},\omega_{L,i}$ & Velocidades de rueda derecha e izquierda del robot $i$ [rad/s] \\
$r_w,b$                 & Radio de rueda y separación entre ruedas del robot diferencial [m] \\
$\eta_o$                & Ganancia de repulsión de obstáculos \\
$\varepsilon_o$         & Regularización del término de obstáculos [m] \\
$c_i$                   & Capacidad del robot $i$ (extensión heterogénea) \\
$M_k$                   & Demanda de capacidad de la carga $k$ (extensión heterogénea) \\
$\gamma$                & Factor de seguridad de capacidad acumulada ($\gamma\geq1$) \\
$a,b,p$                 & Semiejes y exponente de superelipse para cargas rectangulares suavizadas \\
$\tau_{\min}$           & Tiempo muerto mínimo entre transmisiones disparadas por eventos [s] \\
$\delta_p,\delta_W,\delta_h$ & Umbrales de evento para posición, wrench y margen de seguridad \\
$z=(x,v)$               & Estado de fase usado en el MFG de segundo orden \\
$q_L,p_L,H_L$           & Configuración, momento y Hamiltoniano de la carga rígida \\
$w_\rho=(F_\rho,\tau_\rho)$ & Wrench macroscópico aplicado por la densidad del enjambre \\
$d_k^W$                 & Wrench deseado de la carga $k$, compuesto por fuerza y torque objetivo \\
$G_{ik}(q)$             & Mapa de contacto que transforma esfuerzo del robot $i$ en wrench sobre la carga $k$ \\
$G_C(q)$                & Matriz de agarre apilada de una coalición para auditar factibilidad de wrench \\
$\rho_{\partial L}$     & Densidad regularizada de robots sobre el contorno de la carga \\
$B_i$                   & Energía o batería disponible del robot $i$ \\
$B_{\mathrm{crit}}$     & Umbral crítico de batería \\
$E_{\mathrm{return}}$   & Energía estimada necesaria para retornar al cargador \\
$h_i^B$                 & Margen de retornabilidad energética del robot $i$ \\
$P_i$                   & Potencia mecánica instantánea aplicada por el robot $i$ \\
$E_{\mathrm{form}}$     & Energía de mantenimiento de formación alrededor de la carga \\
$\mathcal F_{\mathrm{int}}$ & Funcional integrado económico--físico del framework \\
$\mathcal F_{\mathrm{AIF}}$ & Energía libre operacional asociada al error físico de predicción de la carga \\
$\theta_i$              & Fase angular de contacto del robot $i$ alrededor de la carga \\
$r_k e^{\mathrm{i}\psi_k}$ & Parámetro de orden de Kuramoto para la coalición de la tarea $k$ \\
$\alpha_{ij}$           & Fracción de responsabilidad de evasión ORCA asignada al robot $i$ frente a $j$ \\
$w_i$                   & Peso económico de criticidad del robot $i$ por pérdida de potencial \\
\midrule
\nomhead{Funciones y operadores}
$\Phi(z,n)$       & Sigmoide de cardinalidad \\
$\Phi_c(C,M)$     & Sigmoide de déficit de capacidad acumulada \\
$\Psi(d)$         & Descuento espacial gaussiano \\
$T_k(p)$          & Tiempo, distancia geodésica o potencial eikonal desde $p$ hasta la tarea $k$ \\
$G_k(z)$          & Primitiva de $\Phi_k(z)$ usada para construir potenciales \\
$f_{ik}$          & Payoff del robot $i$ sobre la tarea $k$ \\
$F_{ik}$          & Payoff teórico avanzado, usualmente en función de $\Zeff_k$ o $\zeff_k$ \\
$F_i$             & Campo de potencial artificial sobre el robot $i$ \\
$V(\mathbf{x})$   & Función potencial del juego \\
$E(p,y)$          & Energía maestra que acopla preferencias y movimiento \\
$h(Q)$            & Coste mínimo de desplegar capacidad total $Q$ en una flota heterogénea \\
$D(\lambda)$      & Demanda agregada de capacidad al precio sombra $\lambda$ \\
$\rho(x,t)$       & Densidad espacial suavizada de robots \\
$\rho^N(x,t)$     & Medida empírica suavizada de una flota finita de $N$ robots \\
$\rho_{\max}$     & Densidad máxima admisible antes de activar peaje de congestión \\
$K_h$             & Kernel espacial usado para estimar densidad \\
$v^\tau,T_k^\tau$ & Velocidad admisible y eikonal predictivas \\
$U(x,t)$          & Función de valor del límite de campo medio \\
$U(x,v,t)$        & Función de valor inercial en el MFG de segundo orden \\
$\psi$            & Deseabilidad de Hopf--Cole para linealizar la HJB bajo hipótesis de emparejamiento ruido--control \\
$W_2$             & Distancia 2-Wasserstein entre distribuciones de probabilidad \\
$\Pi_{\mathrm{cong}}$ & Peaje virtual de congestión inducido por densidad futura \\
$Q_{\mathrm{int}}$ & Coste integrado de campo medio con contacto, potencia, batería, formación y congestión \\
$L_G,\lambda_2$   & Laplaciano del grafo $G$ y su valor de Fiedler \\
$\nu_{\mathcal L}$ & Término agregado de conmutación, pérdida de enlaces y retardo del Laplaciano variable \\
$E_{\mathrm{phase}}$ & Energía de sincronización de fases de contacto tipo Kuramoto \\
$\mathcal W_{\mathrm{dem}},\mathcal W_i$ & Wrench demandado por la carga y wrench aportado por el robot $i$ \\
$V_{\mathrm{rig}}$ & Energía de rigidez local de la coalición finita \\
$h_{ia}$ & Función de barrera de seguridad del robot $i$ frente a la restricción local $a$ \\
$u_{\mathrm{nom},i},u_{\mathrm{safe},i}$ & Velocidad nominal y velocidad filtrada por CBF--QP del robot $i$ \\
\midrule
\nomhead{Métricas de evaluación}
$\Theta$                    & \emph{Throughput} de tareas completadas [min$^{-1}$] \\
$\bar{T},T_{\mathrm{coal}}$ & Tiempo medio de tarea y de formación de coalición [s] \\
$F$                         & Tasa de factibilidad ($[0,1]$) \\
$D_{\mathrm{tot}}$          & Distancia total recorrida por la flota [m] \\
$D(R)$                      & Ratio de degradación, $\Theta(R)/\Theta(\infty)$ \\
$\Delta J$                  & Brecha de descentralización frente a T5 \\
$\Delta_{\mathrm{chg}}$     & Cambios de asignación por paso \\
$E_{\mathrm{DAC}}$          & Error medio de estimación de ocupación \\
$N_{\mathrm{fail}}$         & Episodios fallidos por escenario \\
$E_{\mathrm{tot}},E_\tau$   & Energía total consumida y esfuerzo acumulado de torque \\
$\Delta H$                  & Salto de energía port-Hamiltoniana en cambios de rol o contacto \\
$\bar e_{\mathrm{form}}$    & Error medio de mantenimiento de formación \\
$C_{\mathrm{ctrl}}$         & Tiempo de cómputo local del controlador por ciclo \\
$S_a(w)$                    & Puntuación normalizada multi-métrica del tratamiento $a$ con pesos $w$ \\
\bottomrule
\end{longtable}
\endgroup
% La nomenclatura usa longtable solo por maquetación; no debe consumir un número
% de tabla. Se reinicia el contador para que la primera tabla del cuerpo sea Tabla 1.
\setcounter{table}{0}
